\chapter{Semantic Segmentation and the U-Net}
\label{ch:seg}

\chapmeta{Intermediate $\to$ Advanced}{$\sim$6\,h (2.5\,h + 3.5\,h)}{LoveDA (torchgeo)}

\begin{objbox}
\begin{itemize}[leftmargin=1.2em,nosep]
  \item Reframe the task from one label per image to one label per pixel.
  \item Understand the encoder--decoder design and the role of skip connections.
  \item Implement a U-Net from scratch and trace its tensor shapes.
  \item Choose and combine segmentation losses (cross-entropy, Dice).
  \item Train and evaluate a U-Net on the LoveDA land-cover dataset.
\end{itemize}
\end{objbox}

\section{From classification to dense prediction}

Classification answers ``what is in this patch?'' with a single label.
\emph{Semantic segmentation} answers ``what is at every pixel?'', producing an
output of shape $(K,H,W)$ --- a per-pixel class distribution at full input
resolution. The challenge is that the classification encoder repeatedly
\emph{downsamples} (for receptive field and efficiency), so we must recover
spatial detail. The encoder--decoder family does this: an encoder contracts to a
compact, semantically rich representation, and a decoder expands back to full
resolution. The decoder upsamples with transposed convolutions or
interpolation+convolution; the key difficulty is that pooling has discarded the
fine localisation information the decoder needs.

\section{The U-Net architecture}

U-Net (Ronneberger et al., 2015) solves this with \emph{skip connections} that
carry high-resolution encoder feature maps directly to the matching decoder
stage, concatenating them before each decoder convolution
(Figure~\ref{fig:unet}). The encoder provides ``what''; the skips restore
``where''. This single idea is why U-Net dominates medical and EO segmentation.

\begin{figure}[h]
\centering
\begin{tikzpicture}[font=\scriptsize,>=Latex,node distance=4mm]
  \tikzstyle{enc}=[draw,fill=eoblue!12,minimum width=7mm]
  \tikzstyle{dec}=[draw,fill=eogreen!14,minimum width=7mm]
  \tikzstyle{bott}=[draw,fill=eoearth!18,minimum width=7mm]
  % encoder (heights shrink), placed left going down
  \node[enc,minimum height=22mm] (e1) {};
  \node[below=10mm of e1.south,enc,minimum height=16mm,anchor=north] (e2) {};
  \node[below=10mm of e2.south,enc,minimum height=11mm,anchor=north] (e3) {};
  \node[below=10mm of e3.south,bott,minimum height=7mm,anchor=north] (b) {};
  % decoder (heights grow), to the right
  \node[right=42mm of e3,dec,minimum height=11mm] (d3) {};
  \node[above=10mm of d3.north,dec,minimum height=16mm,anchor=south] (d2) {};
  \node[above=10mm of d2.north,dec,minimum height=22mm,anchor=south] (d1) {};
  \node[right=4mm of d1,draw,fill=eored!18,minimum height=22mm,minimum width=5mm] (out) {};
  % down arrows
  \draw[->] (e1.south) -- (e2.north) node[midway,right]{pool};
  \draw[->] (e2.south) -- (e3.north) node[midway,right]{pool};
  \draw[->] (e3.south) -- (b.north) node[midway,right]{pool};
  % up arrows
  \draw[->] (b.east) -| ($(b.east)!0.5!(d3.west)$) |- (d3.west) node[pos=0.9,above]{up};
  \draw[->] (d3.north) -- (d2.south) node[midway,right]{up};
  \draw[->] (d2.north) -- (d1.south) node[midway,right]{up};
  \draw[->] (d1.east) -- (out.west) node[midway,above]{$1\times1$};
  % skip connections
  \draw[->,eored,thick,dashed] (e1.east) -- (d1.west) node[midway,above,black]{skip (concat)};
  \draw[->,eored,thick,dashed] (e2.east) -- (d2.west);
  \draw[->,eored,thick,dashed] (e3.east) -- (d3.west);
  \node[above=1mm of e1]{encoder}; \node[above=1mm of d1]{decoder};
  \node[below=1mm of b]{bottleneck};
  \node[right=0.5mm of out]{\rotatebox{-90}{$K{\times}H{\times}W$ logits}};
\end{tikzpicture}
\caption{U-Net. The encoder (blue) halves resolution and doubles channels at
each stage; the decoder (green) mirrors it, upsampling and halving channels.
Dashed red skip connections concatenate same-resolution encoder features into
the decoder, restoring spatial detail lost to pooling. A final $1\times1$
convolution maps to $K$ class logits per pixel.}
\label{fig:unet}
\end{figure}

A minimal from-scratch U-Net uses a reusable double-convolution block:
\begin{lstlisting}
import torch, torch.nn as nn
class DoubleConv(nn.Module):
    def __init__(self, cin, cout):
        super().__init__()
        self.net = nn.Sequential(
            nn.Conv2d(cin, cout, 3, padding=1, bias=False),
            nn.BatchNorm2d(cout), nn.ReLU(inplace=True),
            nn.Conv2d(cout, cout, 3, padding=1, bias=False),
            nn.BatchNorm2d(cout), nn.ReLU(inplace=True))
    def forward(self, x): return self.net(x)

class Up(nn.Module):
    def __init__(self, cin, cout):
        super().__init__()
        self.up = nn.ConvTranspose2d(cin, cin // 2, 2, stride=2)
        self.conv = DoubleConv(cin, cout)            # cin = up + skip channels
    def forward(self, x, skip):
        x = self.up(x)
        x = torch.cat([skip, x], dim=1)              # concatenate skip
        return self.conv(x)
\end{lstlisting}
Tracing shapes for a $256\times256$ input with base width 64: the encoder
produces feature maps at $256,128,64,32$ resolutions with $64,128,256,512$
channels; the decoder reverses this, each \texttt{Up} consuming the matching
skip; the head outputs $(K,256,256)$.

\section{Loss functions for segmentation}

Per-pixel cross-entropy (Eq.~\ref{eq:ce}, summed over pixels) is the baseline,
but land-cover classes are imbalanced (water and bare soil are rare), and CE is
dominated by frequent classes. The \emph{Dice loss} directly optimises overlap
and is robust to imbalance. For predicted soft probabilities $p$ and one-hot
target $g$ over $N$ pixels and class $k$:
\begin{equation}
\mathcal{L}_{\mathrm{Dice}} = 1 - \frac{1}{K}\sum_{k=1}^{K}
\frac{2\sum_i p_{ik}\,g_{ik} + \varepsilon}{\sum_i p_{ik} + \sum_i g_{ik} + \varepsilon}.
\label{eq:dice}
\end{equation}
In practice a \emph{combined} loss works best, pairing CE's well-behaved
gradients with Dice's overlap objective:
\begin{equation}
\mathcal{L} = \mathcal{L}_{\mathrm{CE}} + \lambda\,\mathcal{L}_{\mathrm{Dice}},
\qquad \lambda \approx 1.
\end{equation}
For severe imbalance, class-weighted CE or focal loss (down-weighting easy
pixels) are alternatives. We measure success with IoU/mIoU, defined in
Chapter~\ref{ch:eval}.

\section{Training on LoveDA}

LoveDA (Wang et al., 2021) provides 5{,}987 high-resolution
($1024\times1024$) scenes over urban and rural domains with seven land-cover
classes, auto-downloaded via \texttt{torchgeo}. The recipe combines everything
so far: tile to $512\times512$ (Chapter~\ref{ch:preproc}), augment image+mask
jointly (Chapter~\ref{ch:aug}), use a U-Net (optionally with a pretrained
encoder, Chapter~\ref{ch:transfer}), optimise the combined loss, and checkpoint
on validation mIoU. The ablation that most convinces students: train the U-Net
\emph{with} and \emph{without} skip connections --- the no-skip variant produces
blurry, mislocalised masks and markedly lower mIoU, making the role of the skips
concrete.

\begin{keybox}
\begin{itemize}[leftmargin=1.2em,nosep]
  \item Segmentation outputs $(K,H,W)$: a class distribution per pixel.
  \item Encoder--decoder recovers resolution; skip connections restore localisation (``what'' + ``where'').
  \item Decoder upsamples via transposed conv (or interpolate+conv) and concatenates skips.
  \item Combine cross-entropy with Dice (Eq.~\ref{eq:dice}) for imbalanced land cover.
  \item Pretrained encoders and joint image+mask augmentation carry straight over from earlier chapters.
\end{itemize}
\end{keybox}

\begin{practbox}
Notebook: \nbpath{chapter\_09\_semantic\_segmentation/09\_practical\_unet\_segmentation.ipynb}.
\begin{enumerate}[leftmargin=1.4em,nosep]
  \item Implement \texttt{DoubleConv}, \texttt{Down}, \texttt{Up}, \texttt{OutConv}
  and assemble a U-Net; print feature-map shapes at every stage.
  \item Implement Dice loss and the combined CE+Dice loss.
  \item Load LoveDA via \texttt{torchgeo}; build a segmentation \texttt{DataLoader}
  with joint augmentation; visualise image/mask pairs.
  \item Train for several epochs; plot loss and validation mIoU curves.
  \item Load the best model and display RGB\,$|$\,GT\,$|$\,prediction triplets.
\end{enumerate}
\end{practbox}

\begin{exbox}
\textbf{9.1 No-skip ablation.} Remove the skip connections; quantify the mIoU
drop and show the blurrier masks.\\[2pt]
\textbf{9.2 Loss comparison.} Train with CE only, Dice only, and CE+Dice; compare
per-class IoU, especially on rare classes.\\[2pt]
\textbf{9.3 Depth/width.} Vary U-Net depth and base width; relate capacity to
mIoU and memory.\\[2pt]
\textbf{Mini-project.} Add an ImageNet-pretrained ResNet encoder to your U-Net
(manual or via SMP, Chapter~\ref{ch:advseg}) and measure the mIoU gain over the
from-scratch version.
\end{exbox}
